\scanpage{159}
\mspage{108}

Supposons que pour \deleted{tout} $t\in \unsure{\dot T|T}\inserted{{}=\bar U}$, le \deleted{\unsure{foncteur}} préfaisceau

\[
\deleted{i_0^{*}(t)\longrightarrow\unsure{\cdot|{}}}
\]

\[
i^{*}(t)=(u\longmapsto\operatorname{Hom}_{X}(u,t))\in U^{\wedge}
\]

soit $\unsure{W|\mathcal W}$-asph. \inserted{(déf. 1 p. 88*)} -- \unsure{p. ex.|i.e.} que la catégorie $U_{/t}$ soit tot. $W$-asph. (prop 8 page 93).

% slanted
\margin{\uline{NB} On dit \unsure{alors} sous ces conditions que $i_0:U\mathbin{-}\bar U$ est \uline{strictement} \deleted{\unsure{tot.}} \uline{$W$-asphérique} (cf déf. 2 page suivante)}

[ NB Le cas où les $U_{/t}$ ont un objet final, cf prop 9 a) i.e. celui où \unsure{l'immersion} $U\to\bar U$ admet un adjoint à droite, est celui envisagé \deleted{\unsure{dans}} chez Dwyer \deleted{\unsure{pour ce}} \deleted{\unsure{avec}} les \uline{imm. ouvertes}

\uline{de Dwyer}\inserted{*})

% slanted
\margin{Je ne connais pas d'autre cas où le th. soit susceptible de s'appliquer, -- les exemples de catégories tot. $W$-asphériques qui n'ont un objet final, \unsure{mais} sont qd même tot. $W$-asphériques, \unsure{$F$|$R$} d'une catégorie $U^{\wedge}$, \unsure{que} je sais sont les $U_{/F}$ où $U$ \unsure{n'importe|quelconque} et $F$ un objet \unsure{filtrant}, ou pour être précis tels que $F$ soit, \unsure{comme} les catégories \uline{filtrantes} sont tot. $W$-asph. si $W$ satisfait \unsure{à des axiomes}, cf démonstration plus loin, en \unsure{des circonstances}, -- p. 112.}

Alors pour tt $f:U\mathbin{-}U'$, l'immersion ouverte

\[
U'\mathbin{-}\bar U'
\]

\unsure{$W$-asphérique} est $=$, \inserted{tout} comme $i_0$, \deleted{\unsure{total}}strictement $W$-asphérique), et a fortiori elle est \unsure{totale|cohér}. \unsure{Par suite|D'ailleurs} (prop. 5) on a ceci

\[
\begin{aligned}
\text{\textcircled{a}}\quad i\in W &\iff i'\in W &&\text{moyennant }W(7\text{ bis})\\
\text{\textcircled{b}}\quad p\in W &\iff q\in W &&\text{moyennant }W(7).
\end{aligned}
\]

\uline{Corollaire 1} \inserted{Supposons $W(7)$.} Si $i:\deleted{\unsure{A}}U\longrightarrow X$ est une imm. ouverte telle que $i_0:U\mathbin{-}\bar U$ soit tot. $W$-asph., \deleted{\unsure{alors}} alors \inserted{$i:U\mathbin{-}X$ est} une $W$-\uline{cofibration} \inserted{(``ultra forte'')}, par quoi j'entends ceci :

pour ttes flèches $U'\xrightarrow{f}U''$ dans $U\backslash\mathrm{Cat}$,
% archive sheet number: 159
